\chapter{Durchfahrtsgesetz und operative Geschwindigkeit}
\label{chap:operational-velocity}

Der Zwei-Geschwindigkeiten-Vergleich in Teil V verwendete konstante
Testgeschwindigkeiten, um die \(v^2\)-Abhängigkeit zu isolieren. Betrieb ist
allgemeiner. Eine Durchfahrt durch das laufende Alignment ist eine monotone
Funktion \(s=s(t)\) in dem Intervall, in dem gewählte Orientierung und
Anwendbarkeitsbedingungen gelten.

\begin{principle}[Geschwindigkeits-Laufzeit-Prinzip]
\label{princ:velocity-runtime}
Geschwindigkeit ist eine qualifizierte Ableitung eines Durchfahrtsgesetzes und
keine intrinsische Alignment-Eigenschaft.
\end{principle}

\section{Raum wird zu Zeit}

\begin{definition}[Geschwindigkeitsbeziehung]
\label{def:velocity-relation}
Für eine differenzierbare Durchfahrt gilt
\[
v(t)=\dot s(t)=\frac{\dd s}{\dd t},
\qquad
a_s(t)=\dot v(t)=\frac{\dd^2s}{\dd t^2}.
\]
Für eine differenzierbare räumliche Größe \(f(s)\) gilt
\[
\frac{\dd}{\dd t}f(s(t))=v(t)f'(s(t)).
\]
\end{definition}

Auf die getrennten Krümmungs- und Überhöhungsgesetze angewendet ergibt sich
\[
\frac{\dd\kappa}{\dd t}=v\kappa'(s),
\qquad
\frac{\dd u}{\dd t}=v u'(s).
\]
Ihre unterschiedlichen Partitionen und Provenienzen bleiben erhalten. Eine
gemeinsame Durchfahrt macht sie weder zu einem Gesetz noch autorisiert sie eine
Kopplungsregel.

Für die zweite Zeitableitung der Krümmung liefert die vollständige Kettenregel
\begin{equation}
\frac{\dd^2\kappa}{\dd t^2}
=v^2\kappa''(s)+\dot v\,\kappa'(s).
\label{eq:curvature-second-time-derivative}
\end{equation}
Der zweite Term entfällt nur unter einer angegebenen
Konstantgeschwindigkeitsannahme. Räumliche Glätte gehört zu \(\kappa(s)\);
zeitliche Anregung hängt zusätzlich von \(s(t)\) ab.

\section{Welche Geschwindigkeit ist gemeint?}

\begin{definition}[Entwurfsgeschwindigkeit]
\label{def:design-velocity}
Eine Entwurfsgeschwindigkeit ist eine Annahme für einen angegebenen
Entwurfszweck mit Quelle, Route oder Elementbereich und Gültigkeit.
\end{definition}

\begin{definition}[Operative Geschwindigkeit]
\label{def:operational-velocity}
Eine operative Geschwindigkeit ist eine geplante, zulässige oder tatsächliche
Durchfahrtseingabe, deren Rolle und Qualifikationen für die Berechnung
angegeben sind.
\end{definition}

Folgende Rollen dürfen nicht zusammenfallen:
\begin{description}
\item[Entwurfsgeschwindigkeit] Annahme zur Entwicklung oder Prüfung eines
Entwurfs.
\item[Zulässige Geschwindigkeit] regel- oder autoritätsgestützte Obergrenze
unter ihrer Anwendbarkeit und Gültigkeit.
\item[Geplante Betriebsgeschwindigkeit] Eingabe aus Fahrplan oder
Betriebsplanung.
\item[Tatsächlich gemessene Geschwindigkeit] zeitindizierte Beobachtung mit
Erfassungsprovenienz und Unsicherheit.
\item[Optimierungsvariable] deklarierter Freiheitsgrad eines Experiments,
kein autorisierter Betriebswert.
\item[Konstante Testgeschwindigkeit] explizite Vereinfachung zur Isolation
einer Abhängigkeit.
\end{description}

\begin{definition}[Geschwindigkeitsumgrenzung]
\label{def:velocity-envelope}
Eine Geschwindigkeitsumgrenzung ist eine qualifizierte Menge zulässiger oder
untersuchter Durchfahrtswerte entlang einer identifizierten Route oder eines
Alignments. Quelle, Bereich, Gültigkeit und Rolle bestimmen ihre Verwendung.
\end{definition}

\section{Kompaktes Beispiel mit veränderlicher Geschwindigkeit}

An derselben Übergangsbogenposition sei
\[
\kappa'=2.0\times10^{-5}\,\mathrm{m^{-2}},
\qquad
\kappa''=-1.0\times10^{-7}\,\mathrm{m^{-3}}.
\]
Bei \(v=20\,\mathrm{m/s}\) gilt
\[
\dot\kappa=4.0\times10^{-4}\,\mathrm{m^{-1}s^{-1}}.
\]
Bei konstanter Geschwindigkeit ist
\(\ddot\kappa=v^2\kappa''=-4.0\times10^{-5}\,
\mathrm{m^{-1}s^{-2}}\). Gilt stattdessen
\(\dot v=0.5\,\mathrm{m/s^2}\), folgt aus
\cref{eq:curvature-second-time-derivative}
\[
\ddot\kappa
=-4.0\times10^{-5}+1.0\times10^{-5}
=-3.0\times10^{-5}\,\mathrm{m^{-1}s^{-2}}.
\]
Das Alignment-Gesetz ist in beiden Fällen identisch. Die Beschleunigung ändert
die zeitliche Exposition durch einen Term, den die
Konstantgeschwindigkeitsvereinfachung auslässt.

\begin{summary}
Jede Verwendung von Geschwindigkeit muss Quelle, Anwendbarkeit, Gültigkeit
und Berechnungsrolle offenlegen. Das Symbol \(v\) allein ist keine
ausreichende Qualifikation.
\end{summary}
